\documentclass[11pt]{article}

\usepackage[margin=1in]{geometry}
\usepackage{amsmath,amssymb,amsthm,mathtools}
\usepackage{microtype}
\usepackage{parskip}
\usepackage{hyperref}

\newtheorem{theorem}{Theorem}[section]
\newtheorem{lemma}{Lemma}[section]
\newtheorem{corollary}{Corollary}[section]
\newtheorem*{remark}{Remark}

\newenvironment{assumptions}{%
  \begingroup
  \renewcommand{\labelenumi}{\textbf{A\arabic{enumi}}.}%
  \renewcommand{\theenumi}{A\arabic{enumi}}%
  \begin{enumerate}
    }{%
  \end{enumerate}
  \endgroup
}

\newcommand{\N}{\mathbb{N}}
\newcommand{\R}{\mathbb{R}}
\newcommand{\1}{\mathbf{1}}
\newcommand{\ind}[1]{\mathbf{1}\{#1\}}
\newcommand{\eqdef}{\coloneqq}
\newcommand{\ens}[1]{\left\{#1\right\}}
\newcommand{\open}[1]{\left(#1\right)}
\newcommand{\closed}[1]{\left[#1\right]}
\newcommand{\modu}[1]{\left|#1\right|}
\newcommand{\mcal}{\mathcal{M}}
\newcommand{\scal}{\mathcal{S}}
\newcommand{\scalt}{\mathcal{S}_{t}}
\newcommand{\bcal}{\mathcal{B}}

\newcommand{\Mcal}{\mcal}
\newcommand{\Scal}{\scal}
\newcommand{\Bcal}{\bcal}
\newcommand{\argmin}{\operatorname*{argmin}}
\newcommand{\argmax}{\operatorname*{argmax}}
\newcommand{\subdiff}{\partial}

\title{Closed-Loop Regret Analysis: Dual Bounds for the All-Ones Mask}
\author{Working Note}
\date{}

\begin{document}
\maketitle

\begin{abstract}
  We study closed-loop replanning for allocating a log-budget under mask constraints, where per-time convex costs are revealed sequentially.
  The main result is a single regret theorem for the replanning policy: regret telescopes into one-step Bellman suboptimalities and is controlled by stage-cost forecast errors and directional errors in the value-function subgradients (dual multipliers).
  We then specialize the statement to convex versus strongly convex costs, to sliding-window masks of width $K$, and to the single all-ones mask.
\end{abstract}

\tableofcontents

%%%%%%%%%%%%%%%%%%%%%%%%%%%%%%%%%%%%%%%%%%%%%%%%%%%%%%%%%%%%%%%%%%%%%%%%%%%%%%%
\section{Setup}
\label{sec:setup}%

Fix a horizon $T \in \N$, a mask family $\mcal_{0} \subseteq \{0,1\}^{T}$, a log-budget $L > 0$, and a threshold cap $\tau_{\max} \in (0,1)$.
Define the log-budget cap $s_{\max} \eqdef -\log({1-\tau_{\max}})$.

\paragraph{Threshold coordinate.}
At each time $t$ the procedure selects a conformal $p$-value threshold $\tau_t \in [0,\tau_{\max}]$.
Given per-time costs $c_1,\dots,c_T$ on $[0,\tau_{\max}]$, the offline design problem is
\begin{equation}
  \label{eq:offline-tau}
  \min_{0 \leq \tau_1,\dots,\tau_T \leq \tau_{\max}} \sum_{t=1}^{T} c_t(\tau_t)
  \quad \text{s.t.} \quad
  \sum_{t=1}^{T} M_t \bigl(-\log({1-\tau_t})\bigr) \leq L, \qquad \forall M \in \mcal_{0}.
\end{equation}

\paragraph{Log transformation.}
Define $s_t \eqdef -\log({1-\tau_t})$, so $s_t \in [0,s_{\max}]$ and $\tau_t = 1-\exp({-s_t})$.
Defining the transformed costs
\begin{equation*}
  \tilde{c}_t(s) \eqdef c_t(1-\exp({-s})),
  \qquad
  s \in [0,s_{\max}],
\end{equation*}
the problem~\eqref{eq:offline-tau} becomes the convex program with linear mask constraints
\begin{equation}
  \label{eq:offline}
  \min_{0 \leq s_1,\dots,s_T \leq s_{\max}} \sum_{t=1}^{T} \tilde{c}_t(s_t)
  \quad \text{s.t.} \quad
  \sum_{t=1}^{T} M_t s_t \leq L, \qquad \forall M \in \mcal_{0}.
\end{equation}

\paragraph{On $s_{\max}$.}
If $\tau_t$ is restricted to the discrete grid $\{0,1/n,\dots,(n-1)/n\}$, then $\tau_{\max}=(n-1)/n$ and $s_{\max}=\log({n})$.
In the remainder we treat $s_{\max}$ as a fixed finite cap.

\paragraph{Closed-loop setting.}
At each time $t$, before observing the realized cost curve $\tilde{c}_t$, the algorithm selects $s_t$ based on the current state $\mathbf{b}_t$ and the forecasted cost curves.
Then $\tilde{c}_t$ is revealed and the incurred cost is $\tilde{c}_t(s_t)$.
Forecasts may be updated over time: we allow the forecast for a future stage $u$ available at time $t$ to differ from the forecast available later.
Concretely, for each $t \in \{1,\dots,T\}$ and each $u \in \{t,\dots,T\}$, let $\hat{c}_{t,u}$ denote the forecast cost curve for stage $u$ available at time $t$.
We treat the realized costs and the entire forecast array $(\hat{c}_{t,u})_{1 \leq t \leq u \leq T}$ as fixed functions; all statements are pathwise.

\paragraph{State and value functions.}
For each mask $M \in \mcal_{0}$, define the remaining budget at time $t$ by
\begin{equation}
  \label{eq:budget-state}
  b_t(M) := L - \sum_{u < t} M_u s_u,
\end{equation}
and let $\mathbf{b}_t := (b_t(M))_{M \in \mcal_{0}} \in \R^{|\mcal_{0}|}$, so $\mathbf{b}_1 = L \1$.
Define the mask-incidence direction at time $t$ as $\mathbf{m}_t := (M_t)_{M \in \mcal_{0}} \in \{0,1\}^{|\mcal_{0}|}$, so the state transition is
\begin{equation}
  \label{eq:state-transition}
  \mathbf{b}_{t+1} = \mathbf{b}_t - s_t \mathbf{m}_t.
\end{equation}

The one-step feasible set is the interval
\begin{equation}
  \label{eq:Scal-def}
  \scalt(\mathbf{b})
  := \left[0, \min\left\{s_{\max}, \inf\{b(M) : M \in \mcal_{0},\ M_t = 1\}\right\}\right],
\end{equation}
with the convention $\inf\emptyset = +\infty$.

For $t \in \{1,\dots,T\}$ and $\mathbf{b} \in \R_{\geq 0}^{|\mcal_{0}|}$, define the remaining offline value function
\begin{equation}
  \label{eq:value-function}
  V_t(\mathbf{b})
  := \min_{\substack{0 \leq s_{t:T} \leq s_{\max} \\
  \sum_{u=t}^{T} M_u s_u \leq b(M),\ \forall M \in \mcal_{0}}}
  \sum_{u=t}^{T} \tilde{c}_u(s_u),
\end{equation}
with terminal condition $V_{T+1}(\cdot) = 0$.
For each planning time $t$, define the time-$t$ forecasted value functions $(\hat{V}_{u \mid t})_{u=t}^{T+1}$ by replacing $\tilde{c}_v$ with $\hat{c}_{t,v}$ in~\eqref{eq:value-function} for stages $v \geq u$, with terminal condition $\hat{V}_{T+1 \mid t}(\cdot)=0$.
In particular, at time $t$ the replanning policy uses $\hat{c}_{t,t}$ and $\hat{V}_{t+1 \mid t}$.

Since the objectives are convex and the constraints are linear, each $V_t$ and each $\hat{V}_{u \mid t}$ is a convex function of $\mathbf{b}$ on $\R_{\geq 0}^{|\mcal_{0}|}$ and is coordinatewise nonincreasing in $\mathbf{b}$.

\paragraph{Standing assumptions.}
\begin{assumptions}
\item\label{ass:A1} (Realized costs.) There exists $\delta>0$ such that, for each $t \in \{1,\dots,T\}$, the realized cost curve $\tilde{c}_t$ is proper, closed, convex, and finite on an open neighborhood $(-\delta,s_{\max}+\delta)$ (hence continuous on $[0,s_{\max}]$).
\item\label{ass:A2} (Forecast costs.) There exists $\delta>0$ such that, for each $t \in \{1,\dots,T\}$ and each $u \in \{t,\dots,T\}$, the forecast curve $\hat{c}_{t,u}$ is proper, closed, convex, and finite on an open neighborhood $(-\delta,s_{\max}+\delta)$ (hence continuous on $[0,s_{\max}]$).
\item\label{ass:A3} (Strong convexity and smoothness.) There exist constants $0 < \mu \leq \beta$ such that, for each $u \in \{1,\dots,T\}$, the realized cost $\tilde{c}_u$ is continuously differentiable on $[0,s_{\max}]$, twice continuously differentiable on $(0,s_{\max})$, and satisfies
  \begin{equation}
    \label{eq:mu-beta-assumptions}
    \mu \leq \tilde{c}_u''(s) \leq \beta
    \qquad \forall s \in (0,s_{\max}).
  \end{equation}
  Moreover, each forecast curve $\hat{c}_{t,u}$ (for all $1 \leq t \leq u \leq T$) is continuously differentiable on $[0,s_{\max}]$, twice continuously differentiable on $(0,s_{\max})$, and satisfies the same bounds $\mu \leq \hat{c}_{t,u}''(s) \leq \beta$ for all $s \in (0,s_{\max})$. In particular, the derivatives $\tilde{c}_u'$ and $\hat{c}_{t,u}'$ are $\beta$-Lipschitz on $(0,s_{\max})$.
\item\label{ass:A4} (All-ones-mask binding and interiority.)
  For the all-ones mask case in Section~\ref{sec:all-ones-mask}, we assume the \emph{binding regime} holds pathwise:
  along the realized trajectory generated by the replanning policy, for every $t \in \{1,\dots,T-1\}$, the remaining budget is strictly positive ($b_{t+1} > 0$) and the tail budget constraint is active for both the true and forecasted tail optima,
  i.e. $\sum_{u=t+1}^{T} s_u^{\star}(b_{t+1})=b_{t+1}$ and $\sum_{u=t+1}^{T} \hat{s}_u(b_{t+1})=b_{t+1}$ (where $\hat{s}_u$ is the forecasted tail allocation produced at time $t$).

  For the average-error bound in Section~\ref{sec:all-ones-mask}, we additionally assume a \emph{pathwise interiority} condition:
  along the realized trajectory, the optimal tail allocations are strictly interior, satisfying $s_u^{\star}(b_{t+1}) \in (0, s_{\max})$ and $\hat{s}_u(b_{t+1}) \in (0, s_{\max})$ for all $u \in \{t+1,\dots,T\}$.
  This naturally holds if the costs are sufficiently homogeneous such that the budget is strictly spread across all available tail stages.
\item\label{ass:A5} (Tail-interiority along the realized trajectory.)
  Along the realized trajectory generated by the replanning policy~\eqref{eq:replan-action}, the state remains tail-interior: for every $t \in \{1,\dots,T\}$ we have $\mathbf{b}_t \in \bcal_t^\circ$.
\end{assumptions}

\paragraph{Replanning policy.}
At each time $t$, the replanning policy selects
\begin{equation}
  \label{eq:replan-action}
  s_t \in \argmin_{s \in \scalt(\mathbf{b}_t)}
  \left[\hat{c}_{t,t}(s) + \hat{V}_{t+1 \mid t}(\mathbf{b}_t - s \mathbf{m}_t)\right],
\end{equation}
and updates $\mathbf{b}_{t+1} = \mathbf{b}_t - s_t \mathbf{m}_t$.
Let $s_t^{\star} \in \argmin_{s \in \scalt(\mathbf{b}_t)} [\tilde{c}_t(s) + V_{t+1}(\mathbf{b}_t - s \mathbf{m}_t)]$ denote any Bellman-optimal action for the true costs.

\paragraph{A raw derivative error metric.}
Under Assumption~\ref{ass:A3}, define the raw derivative forecast error
\begin{equation}
  \label{eq:raw-error}
  e_{t,u} \eqdef \sup_{s \in [0,s_{\max}]} |\tilde{c}_u'(s) - \hat{c}_{t,u}'(s)|,
  \qquad 1 \leq t \leq u \leq T.
\end{equation}
Write $\kappa \eqdef \beta/\mu$.

\subsection{Duality and dual variables}
\label{subsec:duality}

We develop the Lagrangian dual theory for the tail programs~\eqref{eq:value-function}.
Crucially, we employ a \emph{restricted dualization}: we only dualize the linear mask constraints, while keeping the box constraints $0 \leq s_u \leq s_{\max}$ within the primal domain.
We first introduce the Lagrangian and the dual problem over this restricted box, then state the regularity conditions under which Slater's condition ensures strong duality.
Finally, we apply generalized Danskin's theorem to characterize the value-function subgradients, and invoke the Karush-Kuhn-Tucker (KKT) conditions to relate dual multipliers to stage-cost derivatives under interiority.

%--- Step 1: Definitions ---------------------------------------------------

\paragraph{Dual cone.}
Fix $t \in \{1,\dots,T\}$.
A mask $M \in \mcal_{0}$ is \emph{active in the tail} if $\sum_{u=t}^{T} M_u > 0$;
otherwise the constraint $\sum_{u=t}^{T} M_u s_u \leq b(M)$ is trivially satisfied and carries no dual variable.
Define the admissible multiplier cone
\begin{equation*}
  \Lambda_t
  \eqdef
  \left\{\boldsymbol{\lambda} \in \R_{\geq 0}^{|\mcal_{0}|}
  : \lambda(M) = 0 \ \text{whenever}\ \sum_{u=t}^{T} M_u = 0\right\}.
\end{equation*}

\paragraph{Tail Lagrangian.}
Let $\mathbf{s} = (s_t, \dots, s_T)^\top$ denote the vector of tail decisions.
Let $\mathbf{M}_{t:T}$ be the matrix with rows $(M_t, \dots, M_T)$ for each $M \in \mcal_{0}$.
For $\mathbf{b} \in \R_{\geq 0}^{|\mcal_{0}|}$ and $\boldsymbol{\lambda} \in \Lambda_t$, the tail Lagrangian is
\begin{equation*}
  \mathcal{L}_t(\mathbf{s},\boldsymbol{\lambda};\mathbf{b})
  \eqdef
  \sum_{u=t}^{T} \tilde{c}_u(s_u)
  + \boldsymbol{\lambda}^\top (\mathbf{M}_{t:T} \mathbf{s} - \mathbf{b}),
  \qquad \mathbf{s} \in [0,s_{\max}]^{T-t+1}.
\end{equation*}

\paragraph{Dual function.}
Define the dual function by minimizing the Lagrangian over the box:
\begin{equation*}
  \mathcal{D}_t(\boldsymbol{\lambda};\mathbf{b})
  \eqdef
  \inf_{s_{t:T}\in[0,s_{\max}]^{T-t+1}}
  \mathcal{L}_t(s_{t:T},\boldsymbol{\lambda};\mathbf{b}).
\end{equation*}
Note that the infimum is taken over the compact box $[0,s_{\max}]^{T-t+1}$; no extended-value conventions are needed.

\paragraph{Dual-optimal set.}
Define the set of dual optima
\begin{equation*}
  D_t(\mathbf{b})
  \eqdef
  \argmax_{\boldsymbol{\lambda} \in \Lambda_t}
  \mathcal{D}_t(\boldsymbol{\lambda};\mathbf{b}).
\end{equation*}

%--- Step 2: Regularity ----------------------------------------------------

\paragraph{Tail-interior budget region and Slater's condition.}
Slater's condition ensures strong duality for a convex problem if there exists a strictly feasible point.
The duality results below require every active mask constraint to carry a strictly positive remaining budget.
Define the \emph{tail-interior budget region}
\begin{equation*}
  \bcal_t^\circ
  \eqdef
  \left\{\mathbf{b} \in [0,L]^{|\mcal_{0}|}
  : b(M) > 0 \ \text{whenever}\ \sum_{u=t}^{T} M_u > 0\right\}.
\end{equation*}
That is, $\mathbf{b} \in \bcal_t^\circ$ if and only if every mask that is active in the tail has strictly positive remaining budget. This definition allows us to gracefully guarantee a strictly feasible point to claim strong duality.

%--- Step 3: Lemma 1 (Existence & Strong Duality) --------------------------

\begin{lemma}[Existence and strong duality]
  \label{lem:slater-attainment}
  Assume Assumption~\ref{ass:A1}.
  Fix $t \in \{1,\dots,T\}$ and $\mathbf{b} \in \bcal_t^\circ$.
  Then:
  \begin{enumerate}
    \item Slater's condition holds for the tail program~\eqref{eq:value-function}: there exists a strictly feasible point in the box bounds where mask constraints are strictly satisfied.
    \item Strong duality holds:
      $V_t(\mathbf{b})
      = \max_{\boldsymbol{\lambda}\in\Lambda_t}
      \mathcal{D}_t(\boldsymbol{\lambda};\mathbf{b})$.
    \item The dual optimum is attained: $D_t(\mathbf{b}) \neq \emptyset$.
  \end{enumerate}

  Under Assumption~\ref{ass:A2}, the same statements hold for each forecasted tail program defining $\hat{V}_{u \mid t}(\mathbf{b})$ whenever $\mathbf{b} \in \bcal_u^\circ$.
\end{lemma}

\begin{proof}
  \emph{Primal attainment.}
  The feasible set is a closed subset of the compact box $[0,s_{\max}]^{T-t+1}$, hence compact.
  The objective is continuous and convex on the box by Assumption~\ref{ass:A1}.
  Therefore a primal minimizer exists.

  \emph{Slater's condition.}
  Since $\mathbf{b} \in \bcal_t^\circ$, every active mask satisfies $b(M)>0$.
  Let
  \begin{equation*}
    \epsilon
    \eqdef
    \min\left\{\frac{s_{\max}}{2},\
      \min_{\substack{M \in \mcal_{0}:\\[2pt] \sum_{u=t}^{T} M_u>0}}
    \frac{b(M)}{2\sum_{u=t}^{T} M_u}\right\}
    > 0,
  \end{equation*}
  and set $s_u \eqdef \epsilon$ for all $u \in \{t,\dots,T\}$.
  Then $s_u \in (0,s_{\max})$ and, for every active mask,
  \begin{equation*}
    \sum_{u=t}^{T} M_u s_u
    = \epsilon \sum_{u=t}^{T} M_u < b(M).
  \end{equation*}
  Hence a strictly feasible point exists.

  \emph{Strong duality and dual attainment.}
  The tail program is a convex program with linear inequality constraints and a compact feasible set; Slater's condition therefore implies strong duality (zero duality gap).
  For each fixed $\mathbf{s}$, the map $\boldsymbol{\lambda} \mapsto \mathcal{L}_t(\mathbf{s},\boldsymbol{\lambda};\mathbf{b})$ is continuous affine. As the pointwise infimum of continuous affine functions over the compact set $[0,s_{\max}]^{T-t+1}$, the dual function $\mathcal{D}_t(\boldsymbol{\lambda};\mathbf{b})$ is a concave, upper semicontinuous (usc) function.
  Moreover, since $b(M)>0$ for every active mask, taking $\mathbf{s} \equiv 0$ implies there exists a constant $C_0$ such that
  $\mathcal{D}_t(\boldsymbol{\lambda};\mathbf{b}) \leq C_0 - \boldsymbol{\lambda}\cdot\mathbf{b}$
  for all $\boldsymbol{\lambda} \in \Lambda_t$. This implies $\mathcal{D}_t(\boldsymbol{\lambda};\mathbf{b}) \to -\infty$ as $\|\boldsymbol{\lambda}\|\to\infty$ in $\Lambda_t$.
  Since $\mathcal{D}_t$ is upper semicontinuous and coercive, it attains its maximum on some compact level set, and the maximizers form a nonempty compact set. \qedhere
\end{proof}

%--- Step 3: KKT Conditions ------------------------------------------------

\paragraph{KKT conditions and interior solutions.}
Under strong duality, any optimal sequence $\mathbf{s}^\star$ for the tail program and an optimal dual multiplier vector $\boldsymbol{\lambda}^\star \in D_t(\mathbf{b})$ must satisfy the Karush-Kuhn-Tucker (KKT) conditions.

Specifically, assuming continuous differentiability of the stage-costs (Assumption~\ref{ass:A3}), the stationarity condition combined with complementary slackness for the box constraints $0 \leq s_u \leq s_{\max}$ implies that if an optimal coordinate $s_u^\star$ lies strictly in the interior of its bounds ($0 < s_u^\star < s_{\max}$), the multipliers for the box bounds must be exactly zero.
In this regime, the relationship simplifies to equating the unconstrained derivative with the active dual multipliers:
\begin{equation*}
  \tilde{c}_u'(s_u^\star) = - \sum_{M \in \mcal_0} M_u \lambda^\star(M).
\end{equation*}
We rely precisely on this interiority simplification when bounding multiplier errors.

%--- Step 4: Lemma 2 (Danskin / Subgradients) ------------------------------

\paragraph{Value-function subgradients via Danskin's theorem.}
Generalized Danskin's theorem characterizes the subdifferential of a value function evaluated as a maximum over a compact set. We apply it natively to relate the subgradients of our remaining value function $V_t$ directly to the established dual-optimal sets representing marginal constraint values.

\begin{lemma}[Dual multipliers and value-function subgradients]
  \label{lem:dual-subgrad}%
  Fix $t \in \{1,\dots,T\}$ and $\mathbf{b} \in \bcal_t^\circ$.
  Then
  \begin{equation*}
    \subdiff V_t(\mathbf{b})
    = \mathrm{conv}\{-\boldsymbol{\lambda} : \boldsymbol{\lambda} \in D_t(\mathbf{b})\}
    = \{-\boldsymbol{\lambda} : \boldsymbol{\lambda} \in D_t(\mathbf{b})\}.
  \end{equation*}
  The same statement holds for each forecasted value function $\hat{V}_{u \mid t}$ on $\bcal_u^\circ$, with $\tilde{c}_v$ replaced by $\hat{c}_{t,v}$ for $v \geq u$.
\end{lemma}

\begin{proof}
  By Lemma~\ref{lem:slater-attainment}, strong duality holds at $\mathbf{b}$ and the dual maximum is attained on the compact set $D_t(\mathbf{b})$.
  For fixed $\boldsymbol{\lambda}$, the map
  $\mathbf{b} \mapsto \mathcal{D}_t(\boldsymbol{\lambda};\mathbf{b})$
  is affine with gradient $-\boldsymbol{\lambda}$.
  Since the maximum of $\mathcal{D}_t(\cdot;\mathbf{b})$ over $\Lambda_t$ is attained on a compact set and the dual function is affine in $\mathbf{b}$, Danskin's theorem directly gives the subdifferential:
  \begin{equation*}
    \subdiff V_t(\mathbf{b})
    = \mathrm{conv}\{-\boldsymbol{\lambda} : \boldsymbol{\lambda} \in D_t(\mathbf{b})\}.
  \end{equation*}
  Since $D_t(\mathbf{b})$ is convex (as the argmax of a concave function over a convex set), its image under $\boldsymbol{\lambda}\mapsto -\boldsymbol{\lambda}$ is convex, so the convex hull reduces to the set itself. \qedhere
\end{proof}

\subsection{Error quantities}
\label{subsec:errors}

Define the stage-cost subgradient error
\begin{equation}
  \label{eq:eta-t}
  \eta_t(s)
  := \sup_{\substack{\hat{g}_c \in \subdiff \hat{c}_{t,t}(s) \\ g_c \in \subdiff \tilde{c}_t(s)}}
  |\hat{g}_c - g_c|,
  \qquad s \in [0,s_{\max}],
\end{equation}
and let $\eta_t := \sup_{s \in [0,s_{\max}]} \eta_t(s)$.
By Assumptions~\ref{ass:A1} and~\ref{ass:A2}, both subdifferentials are nonempty on $[0,s_{\max}]$ and, since each cost is finite on an open neighborhood of $[0,s_{\max}]$, the subdifferentials are uniformly bounded on $[0,s_{\max}]$, hence $\eta_t$ is finite.

For $\mathbf{b} \in \bcal_{t+1}^\circ$, define the forecasted dual-optimal set
$\hat{D}_{t+1 \mid t}(\mathbf{b}) \eqdef \argmax_{\boldsymbol{\lambda} \in \Lambda_{t+1}} \hat{\mathcal{D}}_{t+1 \mid t}(\boldsymbol{\lambda};\mathbf{b})$,
where $\hat{\mathcal{D}}_{t+1 \mid t}$ is the dual function with $\tilde{c}_u$ replaced by $\hat{c}_{t,u}$ for $u \geq t+1$.
By Lemma~\ref{lem:slater-attainment}, both $D_{t+1}(\mathbf{b})$ and $\hat{D}_{t+1 \mid t}(\mathbf{b})$ are nonempty.
Define the directional dual-multiplier error
\begin{equation}
  \label{eq:epsilon-t}
  \epsilon_t(\mathbf{b})
  := \sup_{\substack{\hat{\boldsymbol{\lambda}} \in \hat{D}_{t+1 \mid t}(\mathbf{b})\\ \boldsymbol{\lambda} \in D_{t+1}(\mathbf{b})}}
  \left|(\hat{\boldsymbol{\lambda}}-\boldsymbol{\lambda})\cdot\mathbf{m}_t\right|,
  \qquad \mathbf{b} \in \bcal_{t+1}^\circ,
\end{equation}
and let $\epsilon_t \eqdef \sup_{\mathbf{b} \in \bcal_{t+1}^\circ} \epsilon_t(\mathbf{b})$.

%%%%%%%%%%%%%%%%%%%%%%%%%%%%%%%%%%%%%%%%%%%%%%%%%%%%%%%%%%%%%%%%%%%%%%%%%%%%%%%

\section{A single regret theorem}
\label{sec:regret}%

\begin{theorem}[Closed-loop regret bound]
  \label{thm:main}%
  Assume Assumptions~\ref{ass:A1},~\ref{ass:A2}, and~\ref{ass:A5}.
  The regret of the replanning policy~\eqref{eq:replan-action} satisfies the exact decomposition
  \begin{equation*}
    \mathrm{Regret}
    \eqdef \sum_{t=1}^{T} \tilde{c}_t(s_t) - V_1(L\1)
    = \sum_{t=1}^{T} \delta_t,
  \end{equation*}
  where $\delta_t \eqdef \tilde{c}_t(s_t) + V_{t+1}(\mathbf{b}_{t+1}) - V_t(\mathbf{b}_t)$ is the one-step Bellman suboptimality, and the one-step bound
  \begin{equation*}
    0 \leq \delta_t \leq \left(\epsilon_t(\mathbf{b}_{t+1}) + \eta_t(s_t)\right)|s_t - s_t^{\star}|
    \leq (\epsilon_t + \eta_t)|s_t - s_t^{\star}|.
  \end{equation*}

  Under Assumption~\ref{ass:A3}, $|s_t - s_t^{\star}| \leq (\epsilon_t + \eta_t)/\mu$ and
  \begin{equation*}
    \mathrm{Regret} \leq \frac{1}{2\mu} \sum_{t=1}^{T} (\epsilon_t + \eta_t)^{2}.
  \end{equation*}
\end{theorem}

\begin{proof}
  We first justify the one-step feasible set~\eqref{eq:Scal-def}.
  Fix $t$ and $\mathbf{b} \geq 0$ coordinatewise.
  If $0 \leq s \leq \min\{s_{\max}, \inf\{b(M) : M \in \mcal_{0},\ M_t = 1\}\}$, then for every mask $M$ with $M_t = 1$ we have $s \leq b(M)$ and for every mask with $M_t = 0$ we have $M_t s = 0 \leq b(M)$.
  Choosing the completion $s_{t+1} = \cdots = s_T = 0$ therefore satisfies the remaining constraints.
  Conversely, if $s > \min\{s_{\max}, \inf\{b(M) : M \in \mcal_{0},\ M_t = 1\}\}$, then either $s > s_{\max}$ or there exists $M$ with $M_t = 1$ and $b(M) < s$, and in both cases feasibility fails because for any completion with $s_{t+1:T} \geq 0$,
  \begin{equation*}
    \sum_{u=t}^{T} M_u s_u \geq M_t s = s > b(M).
  \end{equation*}

  \paragraph{Step 1 (Bellman recursion and telescoping).}
  Define the one-step cost-to-go
  \begin{equation*}
    h_t(s;\mathbf{b}) := \tilde{c}_t(s) + V_{t+1}(\mathbf{b} - s \mathbf{m}_t).
  \end{equation*}
  A standard splitting-and-minimizing argument shows that for every $\mathbf{b} \geq 0$,
  \begin{equation}
    \label{eq:bellman}
    V_t(\mathbf{b}) = \min_{s \in \scalt(\mathbf{b})} h_t(s;\mathbf{b}).
  \end{equation}
  By~\eqref{eq:bellman}, $V_t(\mathbf{b}_t) = \min_{s \in \scalt(\mathbf{b}_t)} h_t(s;\mathbf{b}_t) = h_t(s_t^{\star};\mathbf{b}_t)$.
  Since $\mathbf{b}_{t+1} = \mathbf{b}_t - s_t \mathbf{m}_t$, the Bellman suboptimality can be rewritten as
  \begin{equation*}
    \delta_t = h_t(s_t;\mathbf{b}_t) - \min_{s \in \scalt(\mathbf{b}_t)} h_t(s;\mathbf{b}_t) \geq 0.
  \end{equation*}
  Summing the definition of $\delta_t$ from $t=1$ to $T$ yields
  \begin{equation*}
    \sum_{t=1}^{T} \delta_t
    = \sum_{t=1}^{T} \tilde{c}_t(s_t) + \sum_{t=1}^{T}\left[V_{t+1}(\mathbf{b}_{t+1}) - V_t(\mathbf{b}_t)\right]
    = \sum_{t=1}^{T} \tilde{c}_t(s_t) - V_1(L\1),
  \end{equation*}
  since $V_{T+1}(\cdot)=0$ and $\mathbf{b}_1=L\1$.

  \paragraph{Step 2 (one-step bound and strong-convexity refinement).}
  Fix $t$.
  Write $\hat{h}_t(s;\mathbf{b}) := \hat{c}_{t,t}(s) + \hat{V}_{t+1 \mid t}(\mathbf{b} - s \mathbf{m}_t)$ for the forecasted cost-to-go.
  Because $\hat{c}_{t,t}$ is continuous everywhere on the interval (Assumption~\ref{ass:A2}) and the state transition is linear, we can invoke the subdifferential sum rule and the affine chain rule. Evaluated at $s_t$, every $\hat{g} \in \subdiff \hat{h}_t(s_t;\mathbf{b}_t)$ decomposes as
  \begin{equation*}
    \hat{g} = \hat{g}_c - \hat{g}_V \cdot \mathbf{m}_t,
    \qquad
    \hat{g}_c \in \subdiff \hat{c}_{t,t}(s_t),
    \qquad
    \hat{g}_V \in \subdiff \hat{V}_{t+1 \mid t}(\mathbf{b}_{t+1}),
  \end{equation*}
  and likewise $g = g_c - g_V \cdot \mathbf{m}_t \in \subdiff h_t(s_t;\mathbf{b}_t)$ for any $g_c \in \subdiff \tilde{c}_t(s_t)$, $g_V \in \subdiff V_{t+1}(\mathbf{b}_{t+1})$.
  Since $s_t$ minimizes $\hat{h}_t(\cdot;\mathbf{b}_t)$ over the interval $\scalt(\mathbf{b}_t)$, the normal cone optimality condition implies $0 \in \subdiff \hat{h}_t(s_t;\mathbf{b}_t) + \mathcal{N}_{\scalt(\mathbf{b}_t)}(s_t)$.
  Thus, there exists $\hat{g} \in \subdiff \hat{h}_t(s_t;\mathbf{b}_t)$ such that $\hat{g}(s_t^{\star}-s_t) \geq 0$ for any feasible point $s_t^{\star} \in \scalt(\mathbf{b}_t)$, i.e.\ $\hat{g}(s_t - s_t^{\star}) \leq 0$.
  Combining the subgradient inequality $h_t(s_t^{\star};\mathbf{b}_t) \geq h_t(s_t;\mathbf{b}_t) + g(s_t^{\star}-s_t)$ with $\delta_t = h_t(s_t;\mathbf{b}_t) - h_t(s_t^{\star};\mathbf{b}_t)$ gives
  \begin{equation}
    \label{eq:one-step-chain}
    \delta_t
    \leq g(s_t - s_t^{\star})
    = \underbrace{(g - \hat{g})(s_t - s_t^{\star})}_{\leq |g-\hat{g}| \cdot |s_t - s_t^{\star}|}
    + \underbrace{\hat{g}(s_t - s_t^{\star})}_{\leq 0}
    \leq \bigl(\eta_t(s_t) + \epsilon_t(\mathbf{b}_{t+1})\bigr)|s_t - s_t^{\star}|,
  \end{equation}
  where the last inequality uses $|g-\hat{g}| \leq |\hat{g}_c-g_c| + |(\hat{g}_V-g_V)\cdot\mathbf{m}_t|$ and, by Assumption~\ref{ass:A5}, $\mathbf{b}_{t+1} \in \bcal_{t+1}^\circ$, Lemma~\ref{lem:dual-subgrad} implies $\subdiff V_{t+1}(\mathbf{b}_{t+1})=-D_{t+1}(\mathbf{b}_{t+1})$ and $\subdiff \hat{V}_{t+1 \mid t}(\mathbf{b}_{t+1})=-\hat{D}_{t+1 \mid t}(\mathbf{b}_{t+1})$, so $|(\hat{g}_V-g_V)\cdot\mathbf{m}_t| \leq \epsilon_t(\mathbf{b}_{t+1})$.

  Now assume each $\tilde{c}_t$ is $\mu$-strongly convex on $[0,s_{\max}]$.
  Then $h_t(\cdot;\mathbf{b}_t)$ is $\mu$-strongly convex on $\scalt(\mathbf{b}_t)$.
  By $\mu$-strong convexity, for any $g \in \subdiff h_t(s_t;\mathbf{b}_t)$,
  \begin{equation*}
    h_t(s_t^{\star};\mathbf{b}_t)
    \geq
    h_t(s_t;\mathbf{b}_t) + g(s_t^{\star}-s_t) + \frac{\mu}{2}(s_t-s_t^{\star})^2,
  \end{equation*}
  hence
  \begin{equation*}
    \delta_t
    = h_t(s_t;\mathbf{b}_t)-h_t(s_t^{\star};\mathbf{b}_t)
    \leq g(s_t-s_t^{\star}) - \frac{\mu}{2}(s_t-s_t^{\star})^2.
  \end{equation*}
  Using $g=\hat{g}+(g-\hat{g})$ and $\hat{g}(s_t-s_t^{\star}) \leq 0$ as above gives
  \begin{equation*}
    \delta_t
    \leq |g-\hat{g}| \cdot |s_t-s_t^{\star}| - \frac{\mu}{2}|s_t-s_t^{\star}|^2.
  \end{equation*}
  Let $A_t \eqdef |g-\hat{g}| \leq \eta_t(s_t)+\epsilon_t(\mathbf{b}_{t+1})$.
  Maximizing the quadratic upper bound over $d \eqdef |s_t-s_t^{\star}| \geq 0$ yields
  \begin{equation*}
    \delta_t \leq \frac{A_t^2}{2\mu}
    \leq \frac{1}{2\mu}\bigl(\eta_t(s_t)+\epsilon_t(\mathbf{b}_{t+1})\bigr)^2
    \leq \frac{1}{2\mu}(\epsilon_t+\eta_t)^2.
  \end{equation*}
  The same maximization gives $|s_t-s_t^{\star}| \leq A_t/\mu \leq (\eta_t(s_t)+\epsilon_t(\mathbf{b}_{t+1}))/\mu$.
  Summing over $t$ and using the telescoping identity yields the stated regret bound.
\end{proof}

\section{The all-ones mask}
\label{sec:all-ones-mask}%

\begin{remark}[The all-ones mask]
  \label{rem:all-ones-mask}%
  When $\mcal_{0} = \{\1\}$ consists of a single mask with $M_t = 1$ for all $t$,
  the constraint in~\eqref{eq:offline} reduces to $\sum_{t=1}^{T} s_t \leq L$, the state is scalar $b_t := L - \sum_{u < t} s_u$, and
  \begin{equation*}
    \scalt(b_t) = [0, \min\{b_t, s_{\max}\}].
  \end{equation*}
  The value functions depend only on the scalar budget $b$, and $\mathbf{m}_t = 1$, so the directional dual-multiplier error~\eqref{eq:epsilon-t} reduces to the scalar subgradient discrepancy
  \begin{equation*}
    \epsilon_t(b) = \sup_{\substack{\hat{g} \in \subdiff \hat{V}_{t+1 \mid t}(b) \\ g \in \subdiff V_{t+1}(b)}} |\hat{g} - g|.
  \end{equation*}
\end{remark}

% Bounding the directional multiplier error.
We work under Assumptions~\ref{ass:A3} and~\ref{ass:A4} with $\mcal_{0}=\{\1\}$.
Fix time $t \in \{1,\dots,T-1\}$ and let $b := b_{t+1} \in (0, L]$ be the remaining scalar budget evaluated on the realized trajectory.
Write $N := T-t$.

The scalar tail programs are
\begin{equation*}
  V_{t+1}(b)
  :=
  \min_{\substack{0 \leq s_{t+1:T} \leq s_{\max} \\
  \sum_{u=t+1}^{T} s_u \leq b}}
  \sum_{u=t+1}^{T} \tilde{c}_u(s_u),
  \qquad
  \hat{V}_{t+1 \mid t}(b)
  :=
  \min_{\substack{0 \leq s_{t+1:T} \leq s_{\max} \\
  \sum_{u=t+1}^{T} s_u \leq b}}
  \sum_{u=t+1}^{T} \hat{c}_{t,u}(s_u).
\end{equation*}
For $\lambda \geq 0$ and $u \in \{1,\dots,T\}$ define the \emph{coordinate response maps}
\begin{equation*}
  \psi_u(\lambda) := \argmin_{s \in [0,s_{\max}]} \left(\tilde{c}_u(s) + \lambda s\right),
  \qquad
  \hat{\psi}_{t,u}(\lambda) := \argmin_{s \in [0,s_{\max}]} \left(\hat{c}_{t,u}(s) + \lambda s\right),
\end{equation*}
which are well-defined singletons by strict convexity.
Define the aggregate tail allocations
\begin{equation*}
  \phi(\lambda) := \sum_{u=t+1}^{T} \psi_u(\lambda),
  \qquad
  \varphi(\lambda) := \sum_{u=t+1}^{T} \hat{\psi}_{t,u}(\lambda).
\end{equation*}
Under Assumption~\ref{ass:A3} (strict convexity) and Assumption~\ref{ass:A4} (pathwise interiority), the scalar dual-optimal sets are singletons:
$D_{t+1}(b) = \{\lambda(b)\}$ and $\hat{D}_{t+1 \mid t}(b) = \{\hat{\lambda}(b)\}$.

\begin{theorem}[Multiplier error bound for the all-ones mask]
  \label{thm:epsilon-all-ones}%
  Under Assumptions~\ref{ass:A1}--\ref{ass:A3} and~\ref{ass:A4} with $\mcal_{0}=\{\1\}$ and for $t \in \{1,\dots,T-1\}$,
  \begin{equation}
    \label{eq:epsilon-all-ones}
    \epsilon_t(b)
    = |\lambda(b)-\hat{\lambda}(b)|
    \leq \kappa \bar{e}_{t,>t},
    \qquad
    \bar{e}_{t,>t} := \frac{1}{T-t}\sum_{u=t+1}^{T} e_{t,u}.
  \end{equation}
\end{theorem}

\begin{proof}[Proof (all-ones case)]
  Fix $t \in \{1,\dots,T-1\}$ and write $N := T-t$.

  \paragraph{Step 1 (strong duality, uniqueness, and scalar subgradients).}
  Since $b>0$, the point $s_u := \min\{s_{\max}/2,\,b/(2N)\}$ is strictly feasible for both tail problems, so Slater's condition holds.
  Therefore strong duality holds and the dual maxima are attained for both $V_{t+1}(b)$ and $\hat{V}_{t+1 \mid t}(b)$.
  Under Assumption~\ref{ass:A3}, each tail objective is strictly convex and the feasible set is compact, so the primal optima are unique vectors $s^{\star}_{t+1:T}$ and $\hat{s}_{t+1:T}$.
  Under Assumption~\ref{ass:A4} (interiority and binding), the budget equality constraint binds while the box constraints are strictly loose. This renders the active constraints linearly independent, meaning the corresponding dual multipliers $\lambda(b)$ and $\hat{\lambda}(b)$ are uniquely determined singletons.

  For the scalar value functions, Danskin's theorem yields
  \begin{equation*}
    \partial V_{t+1}(b) = \{-\lambda(b)\},
    \qquad
    \partial \hat{V}_{t+1\mid t}(b) = \{-\hat{\lambda}(b)\},
  \end{equation*}
  hence $\epsilon_t(b)=|\lambda(b)-\hat{\lambda}(b)|$.

  \paragraph{Step 2 (KKT conditions and Mean Value Theorem).}
  The true optimal tail allocation $s^{\star}$ solves the constrained convex program $\min \sum_{u=t+1}^T \tilde{c}_u(s_u)$ subject to $\sum_{u=t+1}^T s_u = b$ and $0 \leq s_u \leq s_{\max}$.
  By Assumption~\ref{ass:A4} (pathwise interiority), $s_u^{\star}$ lies strictly in the open interval $(0, s_{\max})$ for all $u \in \{t+1,\dots,T\}$.
  Thus, by the Karush-Kuhn-Tucker (KKT) conditions and complementary slackness for the box bounds, the boundary multipliers vanish.
  The first-order stationarity condition then holds with exact equality against the unique dual multiplier $\lambda(b) \in \R$ for the budget equality constraint:
  \begin{equation}
    \label{eq:kkt_simplified}
    \tilde{c}_u'(s_u^{\star}) = -\lambda(b).
  \end{equation}
  An identical KKT argument applies to the forecasted optimum $\hat{s}_u \in (0, s_{\max})$ with its unique dual multiplier $\hat{\lambda}(b)$, yielding $\hat{c}_{t,u}'(\hat{s}_u) = -\hat{\lambda}(b)$.
  Define the local partial derivative error $\delta_u := \tilde{c}_u'(\hat{s}_u) - \hat{c}_{t,u}'(\hat{s}_u)$. Notice that $|\delta_u| \leq e_{t,u}$.
  We rewrite the true cost gradient at the forecasted optimum $\hat{s}_u$ as:
  \begin{equation*}
    \tilde{c}_u'(\hat{s}_u) = \hat{c}_{t,u}'(\hat{s}_u) + \delta_u = -\hat{\lambda}(b) + \delta_u.
  \end{equation*}

  By Assumption~\ref{ass:A3}, $\tilde{c}_u'$ is continuously differentiable on the open interval $(0, s_{\max})$. Because both evaluation points lie within this interval, we can directly apply the Mean Value Theorem between $s_u^{\star}$ and $\hat{s}_u$.
  There exists some $\xi_u$ between them such that
  \begin{equation*}
    \tilde{c}_u'(\hat{s}_u) - \tilde{c}_u'(s_u^{\star}) = \tilde{c}_u''(\xi_u)(\hat{s}_u - s_u^{\star}).
  \end{equation*}
  Substituting the derived exact gradients yields:
  \begin{equation*}
    (-\hat{\lambda}(b) + \delta_u) - (-\lambda(b)) = H_u (\hat{s}_u - s_u^{\star}), \qquad \text{where } H_u := \tilde{c}_u''(\xi_u).
  \end{equation*}
  By Assumption~\ref{ass:A3}'s strong convexity blocks, we know $\mu \leq H_u \leq \beta$.
  Rearranging for the coordinate allocation difference gives:
  \begin{equation}
    \label{eq:mvt-coordinate-diff}
    \hat{s}_u - s_u^{\star} = \frac{1}{H_u} (\lambda(b) - \hat{\lambda}(b) + \delta_u).
  \end{equation}

  \paragraph{Step 3 (Aggregating bounded errors).}
  We sum equation~\eqref{eq:mvt-coordinate-diff} over all tail coordinates $u \in \{t+1,\dots,T\}$.
  By the binding regime in Assumption~\ref{ass:A4}, both optima sum identically to the remaining tail budget $b$, enforcing $\sum_u (\hat{s}_u - s_u^{\star}) = b - b = 0$.
  \begin{equation*}
    0 = \sum_{u=t+1}^{T} \frac{1}{H_u} (\lambda(b) - \hat{\lambda}(b)) + \sum_{u=t+1}^{T} \frac{\delta_u}{H_u}.
  \end{equation*}
  Rearranging explicitly for the multiplier difference yields a weighted average of errors:
  \begin{equation*}
    \hat{\lambda}(b) - \lambda(b) = \frac{\sum_{u=t+1}^{T} (\delta_u / H_u)}{\sum_{u=t+1}^{T} (1 / H_u)}.
  \end{equation*}
  Taking absolute values and applying the triangle inequality alongside $|\delta_u| \leq e_{t,u}$, we obtain:
  \begin{equation*}
    |\lambda(b) - \hat{\lambda}(b)| \leq \frac{\sum_{u=t+1}^{T} e_{t,u} / H_u}{\sum_{u=t+1}^{T} 1 / H_u}.
  \end{equation*}
  Since $\frac{1}{\beta} \leq \frac{1}{H_u} \leq \frac{1}{\mu}$, we strictly bound the sums (independently because both sums are strictly positive and the summands are positive scalars):
  \begin{equation*}
    \sum_u \frac{e_{t,u}}{H_u} \leq \frac{1}{\mu} \sum_u e_{t,u}, \qquad \text{and} \qquad \sum_u \frac{1}{H_u} \geq \frac{N}{\beta}.
  \end{equation*}
  This trivially recovers the final bound:
  \begin{equation*}
    |\lambda(b) - \hat{\lambda}(b)| \leq \frac{\frac{1}{\mu} \sum_u e_{t,u}}{\frac{N}{\beta}} = \frac{\beta}{\mu} \left(\frac{1}{N} \sum_{u=t+1}^{T} e_{t,u}\right) = \kappa \bar{e}_{t, >t}.
  \end{equation*}
  \qedhere
\end{proof}

\begin{corollary}[Total regret for the all-ones mask]
  \label{cor:regret-all-ones}%
  Under the assumptions of Theorems~\ref{thm:main} and~\ref{thm:epsilon-all-ones},
  \begin{equation}
    \label{eq:regret-all-ones}
    \mathrm{Regret}
    \leq \frac{1}{2\mu}\sum_{t=1}^{T} (\epsilon_t+e_{t,t})^2
    \leq \frac{\kappa^2}{\mu}\sum_{t=1}^{T-1} \bar{e}_{t,>t}^{2} + \frac{1}{\mu}\sum_{t=1}^{T} e_{t,t}^2.
  \end{equation}
\end{corollary}

\begin{proof}
  Under Assumption~\ref{ass:A3} (differentiability on the closed interval $[0,s_{\max}]$), we have $\subdiff \tilde{c}_t(s_t)=\{\tilde{c}_t'(s_t)\}$ and $\subdiff \hat{c}_{t,t}(s_t)=\{\hat{c}_{t,t}'(s_t)\}$ for all $s_t \in [0,s_{\max}]$, hence $\eta_t(s_t)=|\tilde{c}_t'(s_t)-\hat{c}_{t,t}'(s_t)|\leq e_{t,t}$.
  Therefore Theorem~\ref{thm:main} (strongly convex case) gives $\mathrm{Regret} \leq (2\mu)^{-1}\sum_{t=1}^{T} (\epsilon_t+e_{t,t})^2$ (using $\eta_t(s_t)\le e_{t,t}$).

  In the all-ones case, Theorem~\ref{thm:epsilon-all-ones} gives $\epsilon_t \leq \kappa \bar{e}_{t,>t}$ for $t \in \{1,\dots,T-1\}$ and $\epsilon_T=0$, so
  \begin{equation*}
    \sum_{t=1}^{T} (\epsilon_t+e_{t,t})^2
    \leq 2\sum_{t=1}^{T} \epsilon_t^2 + 2\sum_{t=1}^{T} e_{t,t}^2
    \leq 2\kappa^2\sum_{t=1}^{T-1} \bar{e}_{t,>t}^{2} + 2\sum_{t=1}^{T} e_{t,t}^2,
  \end{equation*}
  so multiplying by $(2\mu)^{-1}$ gives
  \begin{equation*}
    \mathrm{Regret}
    \leq \frac{\kappa^2}{\mu}\sum_{t=1}^{T-1} \bar{e}_{t,>t}^{2} + \frac{1}{\mu}\sum_{t=1}^{T} e_{t,t}^2.
  \end{equation*}
  \qedhere
\end{proof}

% \subsection{Approach 2: Maximum error bound}
% \label{subsec:max-error}

% The average-error bound in Theorem~\ref{thm:epsilon-all-ones} scales with the condition number $\kappa = \beta/\mu$, which can be large if the costs are not very smooth.
% However, we can also prove a robust, absolute maximum-error bound that does not depend on the smoothness constant $\beta$.

% \begin{theorem}[Maximum multiplier error bound for the all-ones mask]
%   \label{thm:epsilon-all-ones-max}%
%   Under Assumptions~\ref{ass:A1}--\ref{ass:A4} with $\mcal_{0}=\{\1\}$ and for $t \in \{1,\dots,T-1\}$,
%   \begin{equation}
%     \label{eq:epsilon-all-ones-max}
%     \epsilon_t(b)
%     = |\lambda(b)-\hat{\lambda}(b)|
%     \leq \max_{u \in \{t+1,\dots,T\}} e_{t,u}.
%   \end{equation}
% \end{theorem}

% \begin{proof}
%   Fix $t \in \{1,\dots,T-1\}$ and let $E := \max_{u \in \{t+1,\dots,T\}} e_{t,u}$.
%   As in the previous proofs, strong duality holds, the primal optima $s^{\star}_{t+1:T}$ and $\hat{s}_{t+1:T}$ are unique, and the dual multipliers $\lambda(b)$ and $\hat{\lambda}(b)$ are unique by the interiority assumption.

%   Assume for contradiction that $\hat{\lambda}(b) < \lambda(b) - E$.
%   Fix any tail coordinate $u \in \{t+1,\dots,T\}$.
%   By the KKT stationarity conditions and Assumption~\ref{ass:A4} (pathwise interiority), both optima are strictly interior, so the bound multipliers vanish. Thus, we have exact equalities:
%   \begin{equation*}
%     \hat{c}_{t,u}'(\hat{s}_u) + \hat{\lambda}(b) = 0 \quad \text{and} \quad \tilde{c}_u'(s_u^{\star}) + \lambda(b) = 0.
%   \end{equation*}
%   By the definition of the uniform derivative error $e_{t,u} \leq E$, we know $\tilde{c}_u'(\hat{s}_u) \geq \hat{c}_{t,u}'(\hat{s}_u) - E$.
%   Combining these yields
%   \begin{equation*}
%     \tilde{c}_u'(\hat{s}_u)
%     \geq -\hat{\lambda}(b) - E
%     > -\lambda(b) + E - E
%     = -\lambda(b) = \tilde{c}_u'(s_u^{\star}),
%   \end{equation*}
%   so $\tilde{c}_u'(\hat{s}_u) > \tilde{c}_u'(s_u^{\star})$.

%   Because $\tilde{c}_u'$ is strictly increasing by Assumption~\ref{ass:A3}, this implies $\hat{s}_u > s_u^{\star}$.
%   Since this strict inequality holds for all $u \in \{t+1,\dots,T\}$, summing over $u$ gives
%   \begin{equation*}
%     \sum_{u=t+1}^{T} \hat{s}_u > \sum_{u=t+1}^{T} s_u^{\star} = b.
%   \end{equation*}
%   This strictly violates the primal binding budget constraint $\sum_u \hat{s}_u = b$, yielding a contradiction.

%   By a symmetric argument (assuming $\lambda(b) < \hat{\lambda}(b) - E$ and swapping the roles of the true and forecasted optima), we also reach a contradiction.
%   Therefore, $|\lambda(b) - \hat{\lambda}(b)| \leq E$, proving the bound.
% \end{proof}

% \begin{corollary}[Total regret via max error]
%   \label{cor:regret-all-ones-max}%
%   Under the assumptions of Theorems~\ref{thm:main} and~\ref{thm:epsilon-all-ones-max},
%   \begin{equation*}
%     \mathrm{Regret}
%     \leq \frac{1}{2\mu}\sum_{t=1}^{T} \left(\max_{u \in \{t+1,\dots,T\}} e_{t,u} + e_{t,t}\right)^2.
%   \end{equation*}
% \end{corollary}

% \begin{proof}
%   This follows exactly as in Corollary~\ref{cor:regret-all-ones}, substituting $\epsilon_t \leq \max_{u > t} e_{t,u}$ from Theorem~\ref{thm:epsilon-all-ones-max}.
% \end{proof}

% %%%%%%%%%%%%%%%%%%%%%%%%%%%%%%%%%%%%%%%%%%%%%%%%%%%%%%%%%%%%%%%%%%%%%%%%%%%%%%%



\section{Window masks of width \texorpdfstring{$K$}{K}}
\label{sec:window-masks}%

Fix $K \in \{1,\dots,T\}$ and define the window family $\mcal_{0}^{\mathrm{win}}(K) \subseteq \{0,1\}^{T}$ by masks $M^{(i)}$, $i \in \{1,\dots,T-K+1\}$, where $M^{(i)}_u = 1$ if $u \in \{i,\dots,i+K-1\}$ and $M^{(i)}_u = 0$ otherwise.
The linear mask constraints become
\begin{equation*}
  \sum_{u=i}^{i+K-1} s_u \leq L, \qquad i \in \{1,\dots,T-K+1\}.
\end{equation*}

At time $t$, let $W_t := \{i \in \{1,\dots,T-K+1\} : t \in \{i,\dots,i+K-1\}\}$ be the set of window constraints that charge $s_t$.
The one-step feasibility bound simplifies to
\begin{equation*}
  \scalt(\mathbf{b}_t) = \left[0, \min\left\{s_{\max}, \min_{i \in W_t}\left(L - \sum_{u=i}^{t-1} s_u\right)\right\}\right],
\end{equation*}
which depends on the past only through the most recent $K-1$ allocations.
Moreover, the directional multiplier error $\epsilon_t$ becomes the discrepancy between subgradients aggregated over the at-most-$K$ active windows that include time $t$.

We introduce a pathwise binding and interiority assumption for window masks, analogous to Assumption~\ref{ass:A4}.
\begin{assumptions}
\item[A6.]\label{ass:A6-win} (Window-mask pathwise binding and interiority.)
  Along the realized trajectory, for every $t \ge K-1$, every tail-active window constraint binds (is a strict equality) for both the true and forecasted tail problems.
  Moreover, the unique optimal actions remain strictly interior: $s_u^{\star} \in (0, s_{\max})$ and $\hat{s}_u \in (0, s_{\max})$ for all $u \in \{t+1,\dots,T\}$.
\end{assumptions}

For window masks, after an initial warm-up period of $K-1$ steps, the directional multiplier error simplifies dramatically.
We characterize the error using a window-residue selector $\gamma_t(u)$, which takes values in $\{0, 1, 2\}$ and periodically selects stages that strongly influence the active window constraints.

\begin{theorem}[Multiplier error bound for window masks]
  \label{thm:window-sensitivity}%
  Under Assumptions~\ref{ass:A1}--\ref{ass:A3} and~A6, for any $t \ge K-1$, the directional multiplier error satisfies
  \begin{equation}
    \label{eq:eps-window-bound-warm}
    \epsilon_t \leq \sum_{u=t+1}^T \gamma_t(u) \tilde{\delta}_{t,u},
  \end{equation}
  where the window-residue selector $\gamma_t(u) \in \{0, 1\}$ is defined as
  \begin{equation}
    \gamma_t(u) := \ind{u \equiv t+1 \pmod{K}}, \label{eq:gamma-window-residue}
  \end{equation}
  and the paired temporal variation error $\tilde{\delta}_{t,u}$ is defined as
  \begin{equation}
    \label{eq:tilde-delta-def}
    \tilde{\delta}_{t,u} :=
    \begin{cases}
      \sup_{b \in \bcal_{t+1}^\circ} \modu{ \left(\tilde{c}_u'(s_u^\star(b)) - \hat{c}_{t,u}'(s_u^\star(b))\right) - \left(\tilde{c}_{u+K-1}'(s_{u+K-1}^\star(b)) - \hat{c}_{t,u+K-1}'(s_{u+K-1}^\star(b))\right) } & \text{if } u+K-1 \le T, \\
      \sup_{b \in \bcal_{t+1}^\circ} \modu{ \tilde{c}_u'(s_u^\star(b)) - \hat{c}_{t,u}'(s_u^\star(b)) } & \text{if } u+K-1 > T.
    \end{cases}
  \end{equation}
  By the triangle inequality, $\tilde{\delta}_{t,u}$ can be uniformly bounded by the individual stage errors: $\tilde{\delta}_{t,u} \le e_{t,u} + e_{t,u+K-1}\ind{u+K-1 \le T} \le 2 \max_{v > t} e_{t,v}$.
\end{theorem}

\begin{proof}
  For $t \geq K-1$, there are $N \eqdef T-t$ tail variables.
  Under Assumption~A6, the binding active tail windows number exactly $N$, indexed by $i \in \{t-K+2,\dots,T-K+1\}$. Let $A \in \R^{N \times N}$ be the incidence matrix for these active constraints, where $A_{\ell,j} = 1$ if $s_{t+j}$ is in the $\ell$-th constraint.
  Because $A$ is a square, lower-triangular matrix with ones on the main diagonal, it is invertible.

  Since all $N$ tail-active inequalities bind at the optimum and $A$ is invertible, the optimal tail allocation is uniquely determined by the linear system $A s = b$. Hence, the optimal set collapses to a singleton, and the true and forecasted optimal actions coincide perfectly: $s^\star = \hat{s} = A^{-1}b$.

  Let $\lambda, \hat{\lambda} \in \R_{\ge 0}^N$ be the dual multipliers for these active constraints. Because the optimal actions $s^\star$ lie strictly in the interior $(0, s_{\max})$, the bound multipliers vanish by complementary slackness.
  The KKT stationarity conditions thus reduce to $A^T \lambda = -\nabla \tilde{c}(s^{\star})$ and $A^T \hat{\lambda} = -\nabla \hat{c}_{t}(s^{\star})$. Since $A^T$ is invertible and the gradients are single-valued (Assumption~\ref{ass:A3}), these multipliers are unique.

  Define the local gradient error $\delta \eqdef \nabla \tilde{c}(s^{\star}) - \nabla \hat{c}_{t}(s^{\star})$, which satisfies $|\delta_u| \le e_{t,u}$. We have $A^T(\hat{\lambda} - \lambda) = \delta$.
  The directional error is an absolute difference of unique vectors:
  \begin{equation*}
    \epsilon_t = \modu{ (\hat{\lambda} - \lambda) \cdot m_t } = \modu{ ((A^{T})^{-1}\delta) \cdot m_t } = \modu{ \delta \cdot (A^{-1}m_t) },
  \end{equation*}
  where $m_t \in \{0,1\}^N$ is the restriction of the time-$t$ incidence vector to these $N$ tail-active constraints.
  Let $w \eqdef A^{-1}m_t$. The vector $w \in \R^N$ is defined by the forward-substitution recurrence $A w = m_t$.
  Because $m_t$ has ones for its first $K-1$ entries and zeros thereafter, the recurrence is:
  $\sum_{j=1}^{r} w_j = 1$ for $r < K$, and $\sum_{j=r-K+1}^{r} w_j = 0$ for $r \geq K$.
  Solving explicitly yields $w_r = 1$ if $r \equiv 1 \pmod{K}$, $w_r = -1$ if $r \equiv 0 \pmod{K}$, and $w_r = 0$ otherwise.

  Substituting $w$ into the inner product bundles the sum into consecutive pairs:
  \begin{equation*}
    \delta \cdot w = \sum_{r=1}^{N} \delta_{t+r} w_r = \sum_{\substack{1 \le r \le N \\ r \equiv 1 \pmod{K}}} \left( \delta_{t+r} - \delta_{t+r+K-1} \ind{r+K-1 \le N} \right).
  \end{equation*}
  Taking the absolute value of the total sum and applying the triangle inequality over the pairs yields
  \begin{equation*}
    \epsilon_t(b) \le \sum_{\substack{1 \le r \le N \\ r \equiv 1 \pmod{K}}} \modu{ \delta_{t+r} - \delta_{t+r+K-1} \ind{r+K-1 \le N} }.
  \end{equation*}
  Mapping indices back to $u = t+r$ identifies the condition $u \equiv t+1 \pmod{K}$, which precisely matches the selector $\gamma_t(u)$.
  The paired difference $\modu{\delta_u - \delta_{u+K-1} \ind{u+K-1 \le T}}$ corresponds exactly to the components defining~\eqref{eq:tilde-delta-def}.
  Taking the supremum over all $b \in \bcal_{t+1}^\circ$ bounds the uniform directional error $\epsilon_t \le \sum_{u=t+1}^{T} \gamma_t(u) \tilde{\delta}_{t,u}$.
\end{proof}

\begin{corollary}[Total regret for window masks]
  \label{cor:regret-window}%
  Under the assumptions of Theorems~\ref{thm:main} and~\ref{thm:window-sensitivity}, the total regret accumulated from time $K-1$ onward is bounded by
  \begin{equation}
    \label{eq:regret-window}
    \mathrm{Regret}_{\geq K-1}
    \eqdef \sum_{t=K-1}^{T} \delta_t
    \leq \frac{1}{2\mu}\sum_{t=K-1}^{T} \left( e_{t,t} + \sum_{u=t+1}^{T} \gamma_t(u) \tilde{\delta}_{t,u} \right)^2.
  \end{equation}
  Alternatively, expressed entirely in terms of the uniform stage errors $e_{t,u}$:
  \begin{equation}
    \label{eq:regret-window-e}
    \mathrm{Regret}_{\geq K-1}
    \leq \frac{1}{2\mu}\sum_{t=K-1}^{T} \left( e_{t,t} + \sum_{u=t+1}^{T} \gamma_t(u) \left(e_{t,u} + e_{t,u+K-1}\ind{u+K-1 \le T}\right) \right)^2.
  \end{equation}
\end{corollary}

\begin{proof}
  By Theorem~\ref{thm:main}, the one-step Bellman suboptimality is bounded by $\delta_t \leq \frac{1}{2\mu}(\epsilon_t + \eta_t)^2$.
  Under Assumption~\ref{ass:A3}, since $s_t \in (0, s_{\max})$, the subdifferentials are singletons and $\eta_t(s_t) = |\tilde{c}_t'(s_t) - \hat{c}_{t,t}'(s_t)| \le e_{t,t}$.
  Substituting the multiplier error bound $\epsilon_t \leq \sum_{u=t+1}^{T} \gamma_t(u) \tilde{\delta}_{t,u}$ from Theorem~\ref{thm:window-sensitivity} yields~\eqref{eq:regret-window}.
  The second bound~\eqref{eq:regret-window-e} follows directly by applying the uniform upper bound $\tilde{\delta}_{t,u} \le e_{t,u} + e_{t,u+K-1}\ind{u+K-1 \le T}$ established in Theorem~\ref{thm:window-sensitivity}.
\end{proof}

%%%%%%%%%%%%%%%%%%%%%%%%%%%%%%%%%%%%%%%%%%%%%%%%%%%%%%%%%%%%%%%%%%%%%%%%%%%%%%%

\end{document}
